\subsection{Validación de calidad de señal}
\label{subsec:validacion_calidad_senal}

Una vez validada la adquisición y seleccionado el montaje experimental, se evaluó la calidad de la señal real. Esta comprobación es necesaria porque una captura puede ser correcta desde el punto de vista de adquisición, manteniendo \SI{250}{\hertz}, sin \textit{sample gaps} y sin estados ADS1299 inválidos, pero contener ventanas contaminadas por ruido de red, movimiento de electrodos, parpadeos, actividad muscular o transitorios de contacto.

La validación se realizó sobre la captura \texttt{20260524-122200\_final\_atenuacion\_artefactos\_mixed\_states}, tomada con el montaje final \texttt{ear\_eeg\_ch1\_only}, \texttt{ADS\_DIAGNOSTIC\_MODE=5} y configuración \texttt{bias\_ch1\_only\_loff\_off}. Esta prueba es útil porque contiene varios estados dentro de una misma adquisición: ojos abiertos, ojos cerrados, mandíbula, recuperación y parpadeo/frente. Por tanto, permite comprobar si el sistema distingue tramos relativamente estables de tramos contaminados sin aceptar o rechazar toda la captura como un único bloque.

La calidad se evalúa por ventanas de análisis. Antes de transformar las características espectrales en parámetros musicales, el backend calcula métricas temporales y espectrales que alimentan el \textit{quality gate}:

\begin{center}
\texttt{compute\_quality\_diagnostics() $\rightarrow$ compute\_spectral\_quality() $\rightarrow$ quality gate}
\end{center}

\begin{figure}[!htbp]
    \centering
    \validacioninclude{width=0.98\textwidth}{images/validacion/eeg_ch1_temporal_750uv_estados.png}
    \caption{Captura temporal completa utilizada para validar la calidad de señal. La representación emplea una escala fija de \(\pm\SI{750}{\micro\volt}\), por lo que los valores que superan ese rango se recortan solo visualmente. Las bandas coloreadas indican los distintos estados experimentales presentes en la adquisición.}
    \label{fig:validacion_mixed_states_temporal_global}
\end{figure}

La Figura~\ref{fig:validacion_mixed_states_temporal_global} muestra que la adquisición se mantiene continua, pero que la calidad no es homogénea. Algunos tramos presentan amplitudes contenidas y otros muestran picos impulsivos o incrementos compatibles con artefactos. Esta diferencia justifica que la señal se analice por ventanas y no como una captura indivisible.

Para resumir la calidad se utilizaron varias métricas: RMS, amplitud pico a pico, componente relativa de \SI{50}{\hertz}, fracción de artefactos e índice de calidad espectral. Estas magnitudes no se interpretan de forma aislada, sino como indicadores complementarios de la validez de cada tramo.

\begin{table}[!htbp]
    \centering
    \small
    \caption{Resumen por estados de la captura utilizada para validar el análisis de calidad.}
    \label{tab:validacion_calidad_estados}
    \begin{tabular}{p{0.23\textwidth} c c c c p{0.23\textwidth}}
        \hline
        \textbf{Estado} & \textbf{RMS med.} & \textbf{PTP med.} & \textbf{50 Hz} & \textbf{Calidad} & \textbf{Diagnóstico} \\
        \hline
        Ojos abiertos & 76,76 & 1272 & 0,259 & 0,987 & Usable con artefactos \\
        Ojos cerrados 1 & 80,01 & 1226 & 0,268 & 0,922 & Artefactos moderados \\
        Mandíbula & 69,72 & 1217 & 0,322 & 0,893 & Usable con artefactos \\
        Recuperación 1 & 80,44 & 1275 & 0,226 & 1,000 & Estable \\
        Parpadeo/frente & 94,22 & 1558 & 0,364 & 0,823 & Usable con artefactos \\
        Recuperación 2 & 109,8 & 1701 & 0,346 & 0,856 & Usable con artefactos \\
        Ojos cerrados 2 & 175,4 & 1725 & 0,188 & 0,621 & Artefacto dominante \\
        \hline
    \end{tabular}
\end{table}

La Tabla~\ref{tab:validacion_calidad_estados} muestra que una misma etiqueta experimental no garantiza una calidad constante. El segundo tramo de ojos cerrados presenta una puntuación de calidad claramente inferior al primer tramo, a pesar de tener una condición nominal parecida. Por ello, la decisión de aceptar, atenuar o descartar información debe realizarse sobre ventanas concretas.

\begin{figure}[!htbp]
    \centering
    \begin{minipage}{0.78\textwidth}
        \centering
        \validacioninclude{width=\textwidth}{images/validacion/fig_00_final_capture_rms_timeline.png}
        \small (a) RMS por ventanas
    \end{minipage}

    \vspace{0.20cm}

    \begin{minipage}{0.78\textwidth}
        \centering
        \validacioninclude{width=\textwidth}{images/validacion/fig_00_final_capture_quality_timeline.png}
        \small (b) Índice de calidad por ventanas
    \end{minipage}

    \vspace{0.20cm}

    \begin{minipage}{0.62\textwidth}
        \centering
        \validacioninclude{width=\textwidth}{images/validacion/fig_11_quality_state_distribution.png}
        \small (c) Distribución de estados de calidad
    \end{minipage}
    \caption{Resumen del comportamiento de calidad en la captura de estados mixtos. El RMS permite detectar incrementos de amplitud, el índice de calidad resume la validez de cada ventana y la distribución final muestra la proporción de ventanas utilizables, dudosas o contaminadas.}
    \label{fig:validacion_calidad_resumen}
\end{figure}

La Figura~\ref{fig:validacion_calidad_resumen} concentra la lectura principal de la validación: la señal real no es uniformemente limpia, pero tampoco debe descartarse por completo. El sistema identifica ventanas con distinta fiabilidad y utiliza esa información para evitar que los tramos contaminados influyan de forma brusca en la sonificación.

Para completar la comprobación se revisaron las señales temporales de los estados principales. La Figura~\ref{fig:validacion_estados_temporales_grid} agrupa las capturas ya generadas durante la validación y permite comparar tramos estables con tramos afectados por artefactos.

\begin{figure}[!htbp]
    \centering
    \begin{minipage}{0.48\textwidth}
        \centering
        \validacioninclude{width=\textwidth}{images/validacion/fig_03_state_ojos_abiertos_reposo_timeseries.png}
        \small (a) Ojos abiertos / reposo
    \end{minipage}
    \hfill
    \begin{minipage}{0.48\textwidth}
        \centering
        \validacioninclude{width=\textwidth}{images/validacion/fig_03_state_ojos_cerrados_reposo_1_timeseries.png}
        \small (b) Ojos cerrados / reposo 1
    \end{minipage}

    \vspace{0.25cm}

    \begin{minipage}{0.48\textwidth}
        \centering
        \validacioninclude{width=\textwidth}{images/validacion/fig_03_state_jaw_timeseries.png}
        \small (c) Mandíbula
    \end{minipage}
    \hfill
    \begin{minipage}{0.48\textwidth}
        \centering
        \validacioninclude{width=\textwidth}{images/validacion/fig_03_state_recuperacion_1_timeseries.png}
        \small (d) Recuperación 1
    \end{minipage}

    \vspace{0.25cm}

    \begin{minipage}{0.48\textwidth}
        \centering
        \validacioninclude{width=\textwidth}{images/validacion/fig_03_state_parpadeo_frente_timeseries.png}
        \small (e) Parpadeo / frente
    \end{minipage}
    \hfill
    \begin{minipage}{0.48\textwidth}
        \centering
        \validacioninclude{width=\textwidth}{images/validacion/fig_03_state_ojos_cerrados_reposo_2_timeseries.png}
        \small (f) Ojos cerrados / reposo 2
    \end{minipage}
    \caption{Comparación temporal por estados en la captura de validación de calidad. La agrupación permite diferenciar tramos relativamente estables de tramos contaminados por movimiento, parpadeo, contacto o actividad muscular.}
    \label{fig:validacion_estados_temporales_grid}
\end{figure}

La inspección temporal se complementó con una revisión espectral básica por estados. Esta comprobación no sustituye a la validación DSP de la Subsección~\ref{subsec:validacion_espectral}; su objetivo es confirmar que los tramos con peor calidad temporal también presentan cambios en frecuencia compatibles con contaminación de la señal.

\begin{figure}[!htbp]
    \centering
    \begin{minipage}{0.48\textwidth}
        \centering
        \validacioninclude{width=\textwidth}{images/validacion/fig_03_state_ojos_abiertos_reposo_psd.png}
        \small (a) Ojos abiertos / reposo
    \end{minipage}
    \hfill
    \begin{minipage}{0.48\textwidth}
        \centering
        \validacioninclude{width=\textwidth}{images/validacion/fig_03_state_ojos_cerrados_reposo_1_psd.png}
        \small (b) Ojos cerrados / reposo 1
    \end{minipage}

    \vspace{0.25cm}

    \begin{minipage}{0.48\textwidth}
        \centering
        \validacioninclude{width=\textwidth}{images/validacion/fig_03_state_mandibula_psd.png}
        \small (c) Mandíbula
    \end{minipage}
    \hfill
    \begin{minipage}{0.48\textwidth}
        \centering
        \validacioninclude{width=\textwidth}{images/validacion/fig_03_state_recuperacion_1_psd.png}
        \small (d) Recuperación 1
    \end{minipage}

    \vspace{0.25cm}

    \begin{minipage}{0.48\textwidth}
        \centering
        \validacioninclude{width=\textwidth}{images/validacion/fig_03_state_parpadeo_frente_psd.png}
        \small (e) Parpadeo / frente
    \end{minipage}
    \hfill
    \begin{minipage}{0.48\textwidth}
        \centering
        \validacioninclude{width=\textwidth}{images/validacion/fig_03_state_ojos_cerrados_reposo_2_psd.png}
        \small (f) Ojos cerrados / reposo 2
    \end{minipage}
    \caption{Comparación espectral por estados en la captura de validación de calidad. La figura no se utiliza para extraer conclusiones fisiológicas, sino para comprobar que los estados con peor calidad temporal también presentan cambios espectrales compatibles con artefactos o contaminación de la señal.}
    \label{fig:validacion_estados_psd_grid}
\end{figure}

La comparación entre las Figuras~\ref{fig:validacion_estados_temporales_grid} y~\ref{fig:validacion_estados_psd_grid} confirma que la calidad de una ventana no depende solo de la amplitud instantánea. También importa cómo la contaminación modifica el contenido frecuencial que después se usará para calcular bandas y características de sonificación.

En conclusión, la captura de estados mixtos demuestra que el sistema no interpreta la señal de forma directa. Primero comprueba su fiabilidad mediante métricas de calidad y después permite que el \textit{quality gate} atenúe o descarte ventanas problemáticas antes de alimentar las etapas de extracción espectral y sonificación.
